\textbf{Phần 1. Trắc nghiệm lựa chọn}
\setcounterpageref{ex}{0}

\begin{ex} %[2H14N73]
$1\text{ m}$ bằng bao nhiêu xăng-ti-mét?
\choice
{$10\text{ cm}$}
{$1000\text{ cm}$}
{$1\text{ cm}$}
{\True $100\text{ cm}$}
\loigiai{
$1\text{ m} = 100\text{ cm}$.\\
\textbf{Đáp án của câu hỏi là D}
}
\end{ex}

\begin{ex} %[2H14N73]
$1\text{ km}$ bằng bao nhiêu mét?
\choice
{\True $1\,000\text{ m}$}
{$100\text{ m}$}
{$10\text{ m}$}
{$10\,000\text{ m}$}
\loigiai{
$1\text{ km} = 1000\text{ m}$.\\
\textbf{Đáp án của câu hỏi là A}
}
\end{ex}

\begin{ex} %[2H14H73]
1 ngày có bao nhiêu giờ?
\choice
{$12$ giờ}
{\True $24$ giờ}
{$30$ giờ}
{$60$ giờ}
\loigiai{
1 ngày = 24 giờ.\\
\textbf{Đáp án của câu hỏi là B}
}
\end{ex}

\begin{ex} %[2H14H73]
1 giờ có bao nhiêu phút?
\choice
{$30$ phút}
{$100$ phút}
{$24$ phút}
{\True $60$ phút}
\loigiai{
1 giờ = 60 phút.\\
\textbf{Đáp án của câu hỏi là D}
}
\end{ex}

\begin{ex} %[2H14H73]
Tháng 1 có bao nhiêu ngày?
\choice
{\True $31$ ngày}
{$30$ ngày}
{$28$ ngày}
{$29$ ngày}
\loigiai{
Tháng 1 có 31 ngày.\\
\textbf{Đáp án của câu hỏi là A}
}
\end{ex}

\begin{ex} %[2H14H73]
Đồng hồ chỉ 8 giờ 15 phút, kim phút đang chỉ vào số nào?
\choice
{Số 12}
{Số 6}
{\True Số 3}
{Số 9}
\loigiai{
15 phút tương ứng với kim phút chỉ số 3.\\
\textbf{Đáp án của câu hỏi là C}
}
\end{ex}

\begin{ex} %[2H14V73]
Tính: $35\text{ kg} + 28\text{ kg} = \dots$
\choice
{\True $63\text{ kg}$}
{$62\text{ kg}$}
{$64\text{ kg}$}
{$65\text{ kg}$}
\loigiai{
$35 + 28 = 63\text{ kg}$.\\
\textbf{Đáp án của câu hỏi là A}
}
\end{ex}

\begin{ex} %[2H14V73]
Tính: $50\text{ l} - 18\text{ l} = \dots$
\choice
{$31\text{ l}$}
{\True $32\text{ l}$}
{$33\text{ l}$}
{$34\text{ l}$}
\loigiai{
$50 - 18 = 32\text{ l}$.\\
\textbf{Đáp án của câu hỏi là B}
}
\end{ex}

\begin{ex} %[2H14V73]
Mai có 2 tờ 500 đồng. Mai có tất cả số tiền là
\choice
{$500$ đồng}
{$800$ đồng}
{\True $1\,000$ đồng}
{$700$ đồng}
\loigiai{
$500 \times 2 = 1000$ đồng.\\
\textbf{Đáp án của câu hỏi là C}
}
\end{ex}

\begin{ex} %[2H14V73]
Một cuộn vải dài $45\text{ m}$, người ta đã bán $18\text{ m}$. Cuộn vải còn lại dài là
\choice
{$26\text{ m}$}
{\True $27\text{ m}$}
{$28\text{ m}$}
{$29\text{ m}$}
\loigiai{
$45 - 18 = 27\text{ m}$.\\
\textbf{Đáp án của câu hỏi là B}
}
\end{ex}

\begin{ex} %[2H14V73]
Bình bắt đầu học bài lúc 19 giờ và học xong lúc 20 giờ. Bình đã học bài trong thời gian là
\choice
{$30$ phút}
{$2$ giờ}
{\True $1$ giờ ($60$ phút)}
{$15$ phút}
\loigiai{
$20 - 19 = 1$ giờ.\\
\textbf{Đáp án của câu hỏi là C}
}
\end{ex}

\begin{ex} %[2H14C73]
Nếu ngày 20 tháng 11 là thứ Tư thì ngày 27 tháng 11 cùng năm đó là thứ mấy?
\choice
{Thứ Ba}
{Thứ Năm}
{Thứ Sáu}
{\True Thứ Tư}
\loigiai{
Cách nhau đúng 7 ngày ($27 - 20 = 7$) nên vẫn là ngày Thứ Tư.\\
\textbf{Đáp án của câu hỏi là D}
}
\end{ex}

\textbf{Phần 2. Trắc nghiệm đúng sai}
\setcounterpageref{ex}{0}

\begin{ex} %[2H14H73]
Xét tính đúng sai về các đơn vị đo đại lượng:
\choiceTF
{\True $1\text{ m} = 10\text{ dm} = 100\text{ cm}$.}
{\True 1 ngày = 24 giờ, 1 giờ = 60 phút.}
{Tháng 4 có 31 ngày.}
{\True Tháng 4 có 30 ngày.}
\loigiai{
\begin{itemize}
    \item $1\text{ m} = 10\text{ dm} = 100\text{ cm}$: Đúng.
    \item 1 ngày = 24h, 1h = 60p: Đúng.
    \item Tháng 4 có 30 ngày (không phải 31). Ý 3 Sai, Ý 4 Đúng.
\end{itemize}
\begin{itemize}
    \item a) \textbf{Đúng}.
    \item b) \textbf{Đúng}.
    \item c) \textbf{Sai}.
    \item d) \textbf{Đúng}.
\end{itemize}
}
\end{ex}

\begin{ex} %[2H14V73]
Một can dầu có $20\text{ l}$, rót đều vào 4 chai nhỏ. Xét tính đúng sai:
\choiceTF
{\True Mỗi chai nhỏ có $5\text{ l}$ dầu ($20 : 4 = 5$).}
{\True 3 chai nhỏ có tất cả $15\text{ l}$ dầu ($5 \times 3 = 15$).}
{Nếu rót vào 5 chai thì mỗi chai có $6\text{ l}$.}
{\True Nếu rót vào 5 chai thì mỗi chai có $20 : 5 = 4\text{ l}$.}
\loigiai{
\begin{itemize}
    \item Mỗi chai 5l: Đúng.
    \item 3 chai 15l: Đúng.
    \item $20 : 5 = 4\text{ l} \neq 6\text{ l}$. Ý 3 Sai, Ý 4 Đúng.
\end{itemize}
\begin{itemize}
    \item a) \textbf{Đúng}.
    \item b) \textbf{Đúng}.
    \item c) \textbf{Sai}.
    \item d) \textbf{Đúng}.
\end{itemize}
}
\end{ex}

\textbf{Phần 3. Trả lời ngắn}
\setcounterpageref{ex}{0}

\begin{ex} %[2H14H73]
$1\text{ m}$ bằng bao nhiêu xăng-ti-mét?
\shortans{100}
\loigiai{
\textbf{Đáp số:} $100$\\
100 cm.
}
\end{ex}

\begin{ex} %[2H14H73]
1 giờ có bao nhiêu phút?
\shortans{60}
\loigiai{
\textbf{Đáp số:} $60$\\
60 phút.
}
\end{ex}

\begin{ex} %[2H14V73]
Tính: $45\text{ kg} + 35\text{ kg}$ theo ki-lô-gam.
\shortans{80}
\loigiai{
\textbf{Đáp số:} $80$\\
$45 + 35 = 80\text{ kg}$.
}
\end{ex}

\begin{ex} %[2H14V73]
Tính: $60\text{ l} - 25\text{ l}$ theo lít.
\shortans{35}
\loigiai{
\textbf{Đáp số:} $35$\\
$60 - 25 = 35\text{ l}$.
}
\end{ex}

\begin{ex} %[2H14V73]
Lan có 1 000 đồng, mua bút hết 700 đồng. Hỏi Lan còn lại bao nhiêu đồng?
\shortans{300}
\loigiai{
\textbf{Đáp số:} $300$\\
$1000 - 700 = 300$ đồng.
}
\end{ex}

\begin{ex} %[2H14C73]
Nếu ngày 5 tháng 5 là thứ Hai thì ngày 12 tháng 5 cùng năm đó là thứ mấy? (Viết: thứ Hai)
\shortans{thứ Hai}
\loigiai{
\textbf{Đáp số:} thứ Hai\\
Cách nhau 7 ngày nên là thứ Hai.
}
\end{ex}
